\chapter{Starting from Zero}
\label{ch:zero}
\begin{chapterintro}
\textbf{You will learn to:} use SI units correctly; convert between units; check
formulas with dimensions; handle scientific notation and significant figures; use
trigonometry to resolve vectors; and read the slope and area of a graph. These are
the tools that every later chapter assumes.
\end{chapterintro}

\section{Physical quantities, units, and dimensions}

A \textbf{physical quantity} is a number together with a unit. Writing
``$v=5$'' is meaningless; writing ``$v=5\ \mathrm{m\,s^{-1}}$'' is a statement
about the world. Mechanics uses three base dimensions:

\begin{center}
\renewcommand{\arraystretch}{1.2}
\begin{tabular}{@{}lll@{}}
\toprule
\textbf{Dimension} & \textbf{Symbol} & \textbf{SI base unit} \\
\midrule
length & $\mathrm{L}$ & metre (m) \\
mass & $\mathrm{M}$ & kilogram (kg) \\
time & $\mathrm{T}$ & second (s) \\
\bottomrule
\end{tabular}
\end{center}

Derived quantities are combinations of these. For example
\[
[v]=\mathrm{L\,T^{-1}},\qquad
[a]=\mathrm{L\,T^{-2}},\qquad
[F]=\mathrm{M\,L\,T^{-2}}.
\]
The last one defines the newton: $1\ \mathrm{N}=1\ \mathrm{kg\,m\,s^{-2}}$.

\begin{keyidea}[Dimensional homogeneity]
Both sides of a correct physical equation must have the same dimensions. You may
add or equate quantities only if they share dimensions. You may multiply and
divide any quantities, combining their dimensions.
\end{keyidea}

\begin{warning}[The most common first-year error]
Adding a length to a time, or a speed to an acceleration, is always wrong. If an
equation contains $x+vt$, then $vt$ must be a length; check it before you
continue.
\end{warning}

\subsection{Prefixes and conversions}

\begin{center}
\renewcommand{\arraystretch}{1.15}
\begin{tabular}{@{}llll@{}}
\toprule
\textbf{Prefix} & \textbf{Symbol} & \textbf{Factor} & \textbf{Example} \\
\midrule
giga & G & $10^{9}$ & $1\ \mathrm{Gm}=10^{9}\ \mathrm{m}$ \\
mega & M & $10^{6}$ & $1\ \mathrm{MPa}=10^{6}\ \mathrm{Pa}$ \\
kilo & k & $10^{3}$ & $1\ \mathrm{km}=10^{3}\ \mathrm{m}$ \\
centi & c & $10^{-2}$ & $1\ \mathrm{cm}=10^{-2}\ \mathrm{m}$ \\
milli & m & $10^{-3}$ & $1\ \mathrm{mm}=10^{-3}\ \mathrm{m}$ \\
micro & $\mu$ & $10^{-6}$ & $1\ \mu\mathrm{m}=10^{-6}\ \mathrm{m}$ \\
nano & n & $10^{-9}$ & $1\ \mathrm{nm}=10^{-9}\ \mathrm{m}$ \\
\bottomrule
\end{tabular}
\end{center}

To convert units, multiply by a fraction equal to one. For example,
\[
72\ \mathrm{km\,h^{-1}}
=72\ \frac{\mathrm{km}}{\mathrm{h}}
\times\frac{1000\ \mathrm{m}}{1\ \mathrm{km}}
\times\frac{1\ \mathrm{h}}{3600\ \mathrm{s}}
=20\ \mathrm{m\,s^{-1}}.
\]

\section{Scientific notation and significant figures}

Write very large or very small numbers as $a\times10^{n}$ with $1\le |a|<10$.
In calculations, keep all digits in intermediate steps and round only the final
answer. A reasonable rule for this book is three significant figures.

\begin{method}[Rounding]
If the first discarded digit is less than 5, round down; if it is 5 or more,
round up. For example, $2.347\to2.35$ (three figures) and $0.004562\to0.00456$.
\end{method}

\section{Trigonometry you must own}

For a right triangle with angle $\theta$,
\[
\sin\theta=\frac{\text{opposite}}{\text{hypotenuse}},\qquad
\cos\theta=\frac{\text{adjacent}}{\text{hypotenuse}},\qquad
\tan\theta=\frac{\text{opposite}}{\text{adjacent}}.
\]
Also
\[
\sin^2\theta+\cos^2\theta=1,\qquad
\tan\theta=\frac{\sin\theta}{\cos\theta}.
\]
Angles are measured in radians in calculus: $\pi\ \mathrm{rad}=180^\circ$. The
arc length of a circle of radius $R$ subtending angle $\theta$ is $s=R\theta$.

\begin{center}
\renewcommand{\arraystretch}{1.2}
\begin{tabular}{@{}lccccc@{}}
\toprule
$\theta$ & $0$ & $30^\circ$ & $45^\circ$ & $60^\circ$ & $90^\circ$ \\
\midrule
$\sin\theta$ & $0$ & $1/2$ & $\sqrt2/2$ & $\sqrt3/2$ & $1$ \\
$\cos\theta$ & $1$ & $\sqrt3/2$ & $\sqrt2/2$ & $1/2$ & $0$ \\
\bottomrule
\end{tabular}
\end{center}

\section{Reading graphs: slope and area}

Two graphical operations recur constantly in mechanics.

\begin{itemize}[leftmargin=1.4em,itemsep=3pt]
  \item The \textbf{slope} of a graph of $y$ against $x$ is $\dd y/\dd x$. On a
        straight line $y=mx+b$, the slope is $m$ and the intercept is $b$.
  \item The \textbf{area} under a graph of $y$ against $x$, between two values of
        $x$, is $\int y\,\dd x$. For a triangle this is $\tfrac12\times\text{base}
        \times\text{height}$; for a rectangle it is base $\times$ height.
\end{itemize}

\begin{keyidea}[What slopes and areas mean in mechanics]
On an $x$--$t$ graph, the slope is velocity. On a $v$--$t$ graph, the slope is
acceleration and the area is displacement. On an $a$--$t$ graph, the area is the
change in velocity.
\end{keyidea}

% ===========================================================================
\section{Worked examples}

\begin{worked}[Dimensional check of a kinematic formula]
Check the dimensions of $v^2=v_0^2+2ax$ and deduce the SI unit of the product
$ax$.
\end{worked}

\begin{solution}{1.1}
Each term must be a velocity squared. Indeed
\[
[v^2]=\mathrm{L^2T^{-2}},\qquad
[2ax]=[a][x]=\mathrm{L\,T^{-2}}\cdot\mathrm{L}=\mathrm{L^2T^{-2}}.
\]
The equation is dimensionally homogeneous. The product $ax$ therefore has the
unit of velocity squared, $\mathrm{m^2\,s^{-2}}$.
\end{solution}

\begin{worked}[Resolving a force into components]
A force of magnitude $F=80.0\ \mathrm{N}$ makes an angle of $35.0^\circ$ above the
positive $x$-axis. Find its $x$- and $y$-components.
\end{worked}

\begin{solution}{1.2}
\[
F_x=F\cos\theta=80.0\cos35.0^\circ=65.5\ \mathrm{N},
\]
\[
F_y=F\sin\theta=80.0\sin35.0^\circ=45.9\ \mathrm{N}.
\]
Check: $F_x^2+F_y^2=(65.5)^2+(45.9)^2=6400$, so $F=80.0\ \mathrm{N}$.
\end{solution}

% ===========================================================================
\section{Exercises and Step-by-Step Solutions}

\begin{exercise}{\warm\ Unit conversions}
Convert: (a) $72\ \mathrm{km\,h^{-1}}$ to $\mathrm{m\,s^{-1}}$;
(b) $1.0\ \mathrm{g\,cm^{-3}}$ to $\mathrm{kg\,m^{-3}}$;
(c) $1.0\ \mathrm{kWh}$ to joules.
\end{exercise}

\begin{tabletworkspace}[6cm]
\end{tabletworkspace}

\begin{solution}{1.1}
(a) $72\ \mathrm{km\,h^{-1}}=72\times1000/3600=20.0\ \mathrm{m\,s^{-1}}$.

(b) $1\ \mathrm{g\,cm^{-3}}=10^{-3}\ \mathrm{kg}/(10^{-2}\ \mathrm{m})^3
=10^{-3}/10^{-6}=10^{3}\ \mathrm{kg\,m^{-3}}$.

(c) $1\ \mathrm{kWh}=10^{3}\ \mathrm{W}\times3600\ \mathrm{s}
=3.60\times10^{6}\ \mathrm{J}$.
\end{solution}

\begin{exercise}{\warm\ Dimensional analysis}
(a) Show that $F=ma$ is dimensionally consistent. (b) The force of a spring is
$F=-kx$, where $x$ is a length. Find the dimensions and SI unit of the spring
constant $k$. (c) Find the dimensions of the gravitational constant $G$ in
$F=GMm/r^2$.
\end{exercise}

\begin{tabletworkspace}[6cm]
\end{tabletworkspace}

\begin{solution}{1.2}
(a) $[ma]=\mathrm{M}\cdot\mathrm{L\,T^{-2}}=\mathrm{M\,L\,T^{-2}}=[F]$.

(b) $k=F/x$, so $[k]=\mathrm{M\,L\,T^{-2}}/\mathrm{L}
=\mathrm{M\,T^{-2}}$, unit $\mathrm{kg\,s^{-2}}$ (equivalently $\mathrm{N\,m^{-1}}$).

(c) From $G=Fr^2/(Mm)$,
\[
[G]=\frac{\mathrm{M\,L\,T^{-2}}\cdot\mathrm{L^2}}{\mathrm{M^2}}
=\mathrm{L^3M^{-1}T^{-2}},
\]
unit $\mathrm{m^3\,kg^{-1}\,s^{-2}}$.
\end{solution}

\begin{exercise}{\warm\ Scientific notation}
Compute
\[
\frac{(3.2\times10^{-3})(4.5\times10^{6})}{1.8\times10^{-2}}
\]
and report the result with three significant figures.
\end{exercise}

\begin{tabletworkspace}[6cm]
\end{tabletworkspace}

\begin{solution}{1.3}
\[
\frac{3.2\times4.5}{1.8}\times10^{-3+6+2}
=8.0\times10^{5}.
\]
The result is $8.00\times10^{5}$ to three significant figures.
\end{solution}

\begin{exercise}{\core\ Right-triangle trigonometry}
In a right triangle the hypotenuse is $13.0\ \mathrm{cm}$ and one leg is
$5.00\ \mathrm{cm}$. Find the other leg, the two acute angles, and the values of
$\sin\theta$, $\cos\theta$, and $\tan\theta$ for the angle opposite the $5.00\ \mathrm{cm}$ leg.
\end{exercise}

\begin{tabletworkspace}[6cm]
\end{tabletworkspace}

\begin{solution}{1.4}
The other leg is
\[
b=\sqrt{13.0^2-5.00^2}=\sqrt{169-25}=\sqrt{144}=12.0\ \mathrm{cm}.
\]
For the angle $\theta$ opposite the $5.00\ \mathrm{cm}$ leg,
\[
\sin\theta=\frac{5.00}{13.0}=0.385,\qquad
\cos\theta=\frac{12.0}{13.0}=0.923,\qquad
\tan\theta=\frac{5.00}{12.0}=0.417.
\]
Thus $\theta=\arcsin(0.385)=22.6^\circ$; the other acute angle is
$90^\circ-22.6^\circ=67.4^\circ$.
\end{solution}

\begin{exercise}{\core\ Components and magnitude}
A displacement has components $\Delta x=-3.00\ \mathrm{m}$ and
$\Delta y=4.00\ \mathrm{m}$. Find its magnitude and the angle it makes with the
positive $x$-axis. Draw the vector.
\end{exercise}

\begin{tabletworkspace}[6cm]
\end{tabletworkspace}

\begin{solution}{1.5}
\[
\norm{\Delta\vec r}=\sqrt{(-3.00)^2+(4.00)^2}=5.00\ \mathrm{m}.
\]
The angle measured from the positive $x$-axis is
\[
\theta=\arctan\frac{4.00}{-3.00}.
\]
The point is in the second quadrant, so $\theta=180^\circ-53.1^\circ=127^\circ$.
\end{solution}

\begin{exercise}{\core\ Slope of a graph}
The position of a particle is measured as a function of time:
\[
\begin{array}{c|ccccc}
t\ (\mathrm{s}) & 0 & 2 & 4 & 6 & 8\\
\hline
x\ (\mathrm{m}) & 2.0 & 6.0 & 10.0 & 14.0 & 18.0
\end{array}
\]
(a) Plot $x$ against $t$. (b) Find the slope and intercept. (c) What physical
quantity does the slope represent? (d) Write the equation of the line.
\end{exercise}

\begin{tabletworkspace}[6cm]
\end{tabletworkspace}

\begin{solution}{1.6}
(a) The points lie on a straight line.

(b) Choosing $(0,2.0)$ and $(8.0,18.0)$,
\[
m=\frac{18.0-2.0}{8.0-0}=2.0\ \mathrm{m\,s^{-1}},\qquad b=2.0\ \mathrm{m}.
\]

(c) The slope is the constant velocity, $v=2.0\ \mathrm{m\,s^{-1}}$.

(d) $x(t)=2.0t+2.0$ with $x$ in metres and $t$ in seconds.
\end{solution}

\begin{exercise}{\core\ Area under a graph}
A car starts from rest and its velocity increases linearly to
$12.0\ \mathrm{m\,s^{-1}}$ in $6.00\ \mathrm{s}$. It then travels at constant
velocity for $4.00\ \mathrm{s}$. Draw the $v$--$t$ graph and find the total
distance travelled.
\end{exercise}

\begin{tabletworkspace}[6cm]
\end{tabletworkspace}

\begin{solution}{1.7}
The $v$--$t$ graph consists of a triangle (from $0$ to $6\ \mathrm{s}$) and a
rectangle (from $6$ to $10\ \mathrm{s}$).
\[
d=\underbrace{\tfrac12(6.00)(12.0)}_{\text{triangle}}
+\underbrace{(4.00)(12.0)}_{\text{rectangle}}
=36.0+48.0=84.0\ \mathrm{m}.
\]
\end{solution}

\begin{exercise}{\core\ The pendulum formula}
The period of a simple pendulum is believed to be
$T=2\pi\sqrt{L/g}$, where $L$ is a length and $g$ an acceleration.
(a) Verify the dimensions. (b) Compute $T$ for $L=1.00\ \mathrm{m}$ and
$g=9.81\ \mathrm{m\,s^{-2}}$. (c) What happens to $T$ if $L$ is quadrupled?
\end{exercise}

\begin{tabletworkspace}[6cm]
\end{tabletworkspace}

\begin{solution}{1.8}
(a) $[L/g]=\mathrm{L}/(\mathrm{L\,T^{-2}})=\mathrm{T^2}$, so
$\sqrt{L/g}$ has dimension $\mathrm{T}$, as required.

(b)
\[
T=2\pi\sqrt{\frac{1.00}{9.81}}=2.01\ \mathrm{s}.
\]

(c) $T\propto\sqrt L$. Quadrupling $L$ multiplies $T$ by $\sqrt4=2$, so the
period doubles to $4.01\ \mathrm{s}$.
\end{solution}

\begin{exercise}{\mastery\ Density of a sphere}
A solid sphere has radius $2.00\ \mathrm{cm}$ and mass $250\ \mathrm{g}$.
(a) Find its density in $\mathrm{kg\,m^{-3}}$. (b) Find its weight on Earth.
(c) If the same material forms a sphere of twice the radius, by what factor does
the mass increase?
\end{exercise}

\begin{tabletworkspace}[6cm]
\end{tabletworkspace}

\begin{solution}{1.9}
(a) The volume is
\[
V=\frac43\pi r^3=\frac43\pi(2.00\times10^{-2})^3
=3.35\times10^{-5}\ \mathrm{m^3}.
\]
The mass is $0.250\ \mathrm{kg}$, hence
\[
\rho=\frac{m}{V}=\frac{0.250}{3.35\times10^{-5}}
=7.46\times10^{3}\ \mathrm{kg\,m^{-3}}.
\]

(b) $W=mg=0.250(9.81)=2.45\ \mathrm{N}$.

(c) $m=\rho V\propto r^3$. Doubling $r$ multiplies the mass by $2^3=8$.
\end{solution}

\begin{exercise}{\mastery\ Linearisation of data}
A student suspects a relation $y=Ax^n$. She plots $\log y$ against $\log x$ and
obtains a straight line of slope $3.0$ passing through
$(\log x,\log y)=(0,0.30)$ (logarithms base 10). Find $n$ and $A$.
\end{exercise}

\begin{tabletworkspace}[6cm]
\end{tabletworkspace}

\begin{solution}{1.10}
Taking base-10 logarithms,
\[
\log y=\log A+n\log x.
\]
The slope of the line is $n=3.0$. At $\log x=0$,
$\log y=\log A=0.30$, hence
\[
A=10^{0.30}=2.0.
\]
Therefore $y=2.0\,x^{3}$.
\end{solution}

